\section{The LLM Era: ChatGPT, Claude, and DeepSeek}

\subsection{From Transformers to Language Models}

The two years following ``Attention Is All You Need'' were a period of rapid experimentation.
In 2018, Google introduced \textbf{BERT} (Bidirectional Encoder Representations from Transformers) \cite{devlin2018bert}, which showed that pre-training a Transformer on massive text corpora and then fine-tuning it for specific tasks outperformed every prior approach in language understanding.

That same year, OpenAI introduced \textbf{GPT} (Generative Pre-trained Transformer) \cite{radford2018improving}, taking the opposite approach: instead of predicting missing words within a sentence, train a Transformer to predict the \emph{next} word, left to right, over billions of tokens of text.

This is called \textbf{language modeling}: the model learns to assign probabilities to sequences of words.
Do it at scale, and something remarkable emerges --- a model that has implicitly absorbed grammar, facts, reasoning patterns, and even style from the training data.

\subsection{The Scaling Insight}

GPT-2 (2019) and GPT-3 (2020) \cite{brown2020language} each scaled this approach dramatically.
GPT-3 had 175 billion parameters --- orders of magnitude larger than anything before it.

The results were qualitatively different, not just quantitatively better.
GPT-3 could write coherent essays, answer questions, translate text, write code, and solve arithmetic problems --- all from a single model, with no task-specific fine-tuning.
Researchers called this \textbf{emergent behavior}: capabilities that appeared at scale without being explicitly trained for.

\begin{tcolorbox}[colback=blue!5!white, colframe=blue!40!black, title=What Are Parameters?]
Parameters are the numerical weights inside a neural network --- the millions (or billions) of numbers that get adjusted during training.
They encode everything the model has learned.
More parameters means more capacity to store and reason about information, though simply adding parameters without enough data and compute gains little.
\end{tcolorbox}

\subsection{ChatGPT: When the World Noticed}

GPT-3 impressed researchers but was accessed through a technical API.
The broader world had little direct contact with it.

That changed in November 2022 when OpenAI released \textbf{ChatGPT}.

ChatGPT was GPT-3.5 with an additional training step: \textbf{Reinforcement Learning from Human Feedback} (RLHF) \cite{ouyang2022training}.
Human raters evaluated model responses and ranked them; a separate model learned to predict those rankings; and the language model was then fine-tuned to produce responses that the ranking model would score highly.

The result was a model that felt remarkably natural to talk to.
It remembered the conversation, acknowledged its mistakes, and declined harmful requests.
ChatGPT reached one million users in five days.
It remains one of the fastest product adoptions in history.

\subsection{Claude}

\textbf{Claude}, developed by Anthropic, followed a similar recipe but with a different emphasis.
Anthropic was founded by former OpenAI researchers who wanted to focus more directly on AI safety.

Claude incorporated a technique called \textbf{Constitutional AI} \cite{bai2022constitutional}: rather than relying solely on human raters, the model was trained using a written set of principles (a ``constitution'') that guided its own self-critique and revision of its outputs.
This approach aimed to make the model's values more transparent and consistent.

Claude has become known for producing long, nuanced, carefully reasoned responses --- and for being willing to say ``I'm not sure'' when it is not.

\subsection{DeepSeek}

In early 2025, a Chinese AI lab called \textbf{DeepSeek} released a model that stunned the field.
DeepSeek-R1 \cite{deepseek2025r1} matched or exceeded the performance of models from OpenAI and Google on several reasoning benchmarks --- and did so at a fraction of the training cost.

The team achieved this through several innovations, including a \textbf{Mixture of Experts} (MoE) architecture that activates only a subset of the model's parameters for any given input, and aggressive efficiency optimizations throughout training.

DeepSeek made two moves that sent ripples through the industry:
\begin{enumerate}
    \item It released its model weights openly, so anyone could download and run the model.
    \item It published detailed technical reports explaining how it achieved its results.
\end{enumerate}

This moment reframed the conversation about AI development: frontier performance was no longer the exclusive domain of the largest and best-funded labs.
